% ============================================================================
%  Stacks: Overview and Objectives  (CIS 4020)
%  Requires classnotes.sty in the same folder or your local texmf tree.
%  Compile twice with pdflatex so "Page X of Y" resolves.
% ============================================================================
\documentclass[11pt]{article}
\usepackage[margin=1in]{geometry}
\usepackage[pagenumbers]{classnotes}

\classnotesmeta{Stacks: Overview and Objectives (CIS 4020)}

\begin{document}

\classnotestitle{CIS 4020 \quad$\cdot$\quad DATA STRUCTURES}{Stacks: Overview and Objectives}{An abstract data type built on last in, first out order}

\section{What this unit is about}

A stack is an abstract data type, the same kind of construct as the linked list from the
last unit. As with a linked list, we care about the wall between the operations a stack
offers and how those operations get carried out behind the scenes, and a stack can be
built on more than one foundation. Where the linked list unit compared array-based and
pointer-based implementations of a list, this unit does the same comparison for a stack:
an array-based version and a linked version, built on the node and link ideas you already
know.

A stack holds its items in last in, first out order, often shortened to LIFO. Only one end
of the stack, called the top, is available for adding or removing an item. The two core
operations are push, which adds an item to the top, and pop, which removes and returns the
item at the top. A third operation, peek or top, looks at the top item without removing it.
Stacks show up constantly once you know to look for them: undo history in an editor,
the back button in a browser, matching parentheses or brackets, evaluating an arithmetic
expression, and the call stack that keeps track of recursive calls from the first unit of
this course.

As with the earlier units, we work in \textbf{two languages at once}. Computer Science
students write \textbf{C/C++}; CIS students write \textbf{Java}. The stack idea is
identical in both, and both the array-based and linked implementations appear in each
language.

Class time favors drawings and live code over slides. Read ahead, try tracing a push and
pop sequence on paper yourself, and come ready to think out loud.

\section{Learning objectives}

By the end of the unit you should be able to:

\subsection*{Remember and understand}
\begin{itemize}[leftmargin=1.4em]
  \item State that a stack is an ADT, and recall the wall/information-hiding idea from the
    linked list unit as it applies to a stack.
  \item Define a stack and state its defining property, last in, first out order.
  \item Discuss the standard stack operations: push, pop, peek (or top), isEmpty, and size.
  \item Discuss the two implementations covered in this unit, array-based and linked, and
    recall the node and link vocabulary from the linked list unit as it applies to a linked
    stack.
  \item List common uses of a stack, including undo history, browser back buttons, matching
    parentheses or brackets, expression evaluation, and the recursive call stack.
\end{itemize}

\subsection*{Analyze}
\begin{itemize}[leftmargin=1.4em]
  \item Compare the array-based and linked implementations of a stack, including how each
    handles growth, memory use, and the cost of push and pop.
  \item Explain the relationship between the stack ADT and the run-time call stack used
    when a program executes a recursive function.
  \item Trace a stack by hand through a sequence of push, pop, and peek operations, and
    state its contents and top at each step.
  \item Judge whether a stack is the right ADT for a given problem, as opposed to a plain
    linked list or another structure already covered.
\end{itemize}

\subsection*{Apply}
\begin{itemize}[leftmargin=1.4em]
  \item Diagram a stack by hand, showing the top and the effect of a push or a pop.
  \item Code stack operations from scratch in both an array-based version and a linked
    version, in your language.
  \item Apply a stack to a standard problem, such as checking that parentheses or brackets
    are balanced, or evaluating a postfix expression.
  \item Generate a stack-based code solution to a new problem.
\end{itemize}

\subsection*{Evaluate}
\begin{itemize}[leftmargin=1.4em]
  \item Select the more appropriate stack implementation, array-based or linked, for a
    given scenario, and justify the choice against the alternative.
\end{itemize}

\end{document}
